\begin{theorem}[Wild odd correction preparation]
\label{mod:cases-wildodd-correctionpreparation}
Starting from the source-retaining phase package, construct the actual
\texttt{WildOddCorrectionData} and bridge its correction units to the indexed
presentation required by the High and Low table wrappers.  Prove the
correction-depth memberships and the pointwise character linearizations, then
construct the actual-row \texttt{ExactOddPhaseAssembly}.  Transport the raw
critical quotient to its exact \texttt{residualPhase\_eq}.

After the assembly is fixed, construct its correction coordinate and
valuation depth, choose the canonical nonidentity norm character $\tau$, prove
the general stationary-class \texttt{tauLinearization}, and return the complete
\texttt{WildOddNonstableHigherData} consumed by the existing main theorem.
\LeanFile{LanglandsFirstMainLemma/Cases/WildOdd/CorrectionPreparation.lean}
\lean{LanglandsFirstMainLemma.wildOddNonstableHigherData}
\leanok
\SourceText{Lemmas lem:parameter-ratios, lem:correction-depth, and lem:p-valuation; Corollary cor:wild-odd.}
\uses{mod:cases-wildodd-phasepreparation,mod:cases-wildodd-correctionvaluation,mod:ramification-normcharacters,mod:ramification-traceideals,mod:cases-wildodd-main}
\FormalizationContract
The order is correction datum, depth and pointwise linearizations, actual
table-backed assembly, residual identity, then the correction coordinate and
general $\tau$-linearization.  The correction element is the source-tied
assembly coordinate and every stationary linearization is constructed at its
actual depth.  Expected implementation size: 1200--1800 Lean lines.
\end{theorem}
